\chapter{Classical and Constructive Logic}
\label{ch:classical-constructive}

This chapter makes the constructive/classical boundary a concrete,
observable formal fact rather than historical or terminological
background: one direction of double negation typechecks without
qualification, and the other genuinely requires a classical
principle to be invoked by name.

\section{Double Negation Introduction}

\begin{experiment}{$p \to \lnot\lnot p$}
\begin{lean}
theorem dn_intro (p : Prop) (hp : p) : ¬¬p :=
  fun hnp => hnp hp
\end{lean}
This direction is constructively valid: it requires nothing beyond
the definition of negation as $p \to \mathrm{False}$.
\end{experiment}

\section{Double Negation Elimination}

Try to prove $\lnot\lnot p \to p$ with only the tactics from
Chapter~\ref{ch:natural-deduction} and the attempt stalls
immediately: \texttt{intro h} leaves a hypothesis
$h : \lnot\lnot p$, but none of \texttt{exact}, \texttt{apply}, or
\texttt{cases} can turn a doubly-negated hypothesis into a proof of
$p$ itself — \texttt{cases} has nothing to split, since $\lnot\lnot p$
is not built from a multi-constructor type the way $p \lor q$ or
$\exists x, P\,x$ are. This is not a countermodel — the statement is
true — but the gap is real: nothing introduced so far closes it.

\begin{experiment}{Double negation elimination, classically}
\begin{lean}
theorem dn_elim (p : Prop) (h : ¬¬p) : p :=
  Classical.byContradiction h
\end{lean}
This looks almost too easy, and that is the point: unfolding
notation, $\lnot\lnot p$ \emph{is} $(\lnot p) \to \mathrm{False}$,
and \texttt{Classical.byContradiction} has exactly that as its
hypothesis type. The classical principle is not doing subtle work
here — it is being invoked directly, by name, at precisely the one
step Chapter~\ref{ch:natural-deduction}'s toolkit could not take.
\end{experiment}

\section{Excluded Middle}

\texttt{Classical.em} states $p \lor \lnot p$ for every proposition
$p$, and given the rest of classical logic it is interchangeable
with double negation elimination — each is provable from the other.
One direction was already used implicitly above; the other direction
is worth deriving explicitly, since it reuses \texttt{cases} from
Chapter~\ref{ch:natural-deduction} directly.

\begin{experiment}{Double negation elimination from excluded middle}
\begin{lean}
theorem dn_elim_from_em (p : Prop) (h : ¬¬p) : p := by
  cases Classical.em p with
  | inl hp => exact hp
  | inr hnp => exact absurd hnp h
\end{lean}
\texttt{Classical.em p} supplies $p \lor \lnot p$ directly; the
first case closes immediately, and the second uses
\texttt{absurd}\footnote{\texttt{absurd : a → ¬a → b}, valid for any
target type \texttt{b} — from a proof of \texttt{a} and a proof of
\texttt{¬a}, anything follows.} to extract a contradiction from
$\lnot p$ together with $h : \lnot\lnot p$.
\end{experiment}

Constructive logic does not reject $p \lor \lnot p$ because it is
false — under classical assumptions it plainly is not — but because
treating it as a primitive would mean asserting, for
\emph{every} proposition whatsoever, that one of its two sides is
decidable in principle, with no general procedure supplied for
determining which. A constructive proof of $p \lor \lnot p$ has to
actually produce one side or the other; for a proposition about, say,
an open mathematical question, no such construction is available. In
Lean's classical fragment this worry is set aside by fiat, via the
axiom underlying \texttt{Classical.em} itself.

\section{Exercises}

\begin{exercise}
Prove Peirce's law, $((p \to q) \to p) \to p$, and identify exactly
where a classical principle is required.
\end{exercise}
